% =============================================================================
% Chapter 6: Cross-Sector Integration and the Common Medium
% Input-ready fragment for the lean toy-model research proposal.
% This file intentionally contains no preamble or document environment.
% =============================================================================

\chapter{Cross-Sector Integration and the Common Medium}
\label{ch:cross-sector-integration}

The light and far-field chapters do not define several independent effective
models. They define a linked set of conditions that one parent medium must meet.
This chapter performs the transition from those conditions to an actual
realization attempt. It normalizes the sector outputs in one Medium Capability
Matrix, identifies the quantities that several sectors share, tests whether a
common admissible region exists, and states what must be frozen in Candidate
Parent System V1. The candidate is then required to produce one stable brane and
one complete response spectrum from which the light and far-field results are
rederived.

The compatibility burdens collected here are not established contradictions.
The companion cross-sector review found no present no-go theorem. Each burden
becomes a falsification gate only when a frozen system shows that a mandatory
capability is impossible or necessarily entails a forbidden consequence. The
central risk is therefore not an already-proved inconsistency, but the
possibility that requirements which are separately plausible have an empty
intersection when imposed on one medium.

\section{The Medium Capability Matrix}

The Medium Capability Matrix is the common interface between requirement
discovery and model construction. It is not a narrative summary and it is not a
scorecard. Each row identifies one controlled physical or mathematical object,
the property required of it, the regime in which that property is required, the
sectors that use it, and the evidence needed to decide whether the requirement
is satisfied. Two sectors that refer to the same mode, coefficient, source, or
normalization must point to the same row rather than introducing parallel
versions of it.

The matrix must preserve four different kinds of entry. First are
\emph{candidate choices}: the field inventory, dynamical branch, operator basis,
symmetries, independent coefficients, and boundary or reservoir ensemble.
Second are \emph{background outputs}: the slab profile and thickness, guided
normal profiles, mode compositions, eigenvalues or retarded poles, continuum
thresholds, and physical normalizations. Third are \emph{provisional source
contracts}: the parity, multipole content, continuity law, and symbolic form
factor required of an eventual throat. Fourth are \emph{later nonlinear
outputs}: actual throat source magnitudes, universality, inertial and passive
response, pair topology, and the converged reservoir state. Keeping these
classes separate prevents an unknown throat or global quantity from being
promoted to an adjustable parent parameter.

A compact matrix schema is shown in \cref{tab:capability-matrix-schema}. The
filled matrix remains a working companion to the proposal; the proposal needs
only enough structure to make the common-medium decision unambiguous.

\begin{table}[htbp]
\centering
\caption{Compact schema for the Medium Capability Matrix.}
\label{tab:capability-matrix-schema}
\begin{tabularx}{\textwidth}{@{}L{0.20\textwidth}L{0.34\textwidth}Y@{}}
\toprule
\textbf{Entry class} & \textbf{Controlled content} & \textbf{Cross-sector use} \\
\midrule
Background and branch
& Arena, fields, conservative, relaxational, or mixed dynamics, finite-slab
  profile, shear support, constraints, and reservoir variables
& Establishes the common host on which every mode, source, force, and lifetime
  calculation is performed. \\
Mode or response branch
& Field composition, reflection parity, normal profile, static stiffness,
  dispersion or retarded pole, gap, continuum placement, norm or energy sign,
  damping, leakage, and observable residue
& Connects guided light, the electric odd and even sectors, longitudinal and
  bulk response, radiation, throat support, critical velocity, and drag. \\
Source and coupling
& Odd and even source vectors, form factors, current or constraint identity,
  one-throat coupling magnitude, gravity-source reduction, and mode overlap
& Connects static forces to moving-source response, acceleration radiation,
  universality, and the later replacement of provisional sources by solved
  throats. \\
Observable and ensemble
& Force or flux extraction, equilibrium, periodic, or driven ensemble,
  outgoing conditions, one-time calibration, and declared physical regime
& Prevents a static energy, retarded response, mode normalization, or
  calibration from being changed between sectors. \\
Decision record
& Exact gate, bounded departure, or scope choice; metric and domain where
  needed; uncertainty; current evidence; failure scope; and dependencies
& Converts qualitative words such as ``small,'' ``long lived,'' and
  ``universal'' into a conjunctive common-medium decision. \\
\bottomrule
\end{tabularx}
\end{table}

For every background eigenmode, the matrix should record enough information to
serve both the low-wave-number force audit and the finite-wave-number light
program. A label such as ``longitudinal mode'' is insufficient. The same row
must state its actual mixture of displacement, interface, density, order, bulk,
and reservoir variables; its parity; its profile across the slab; its static
and dynamic behavior; its source overlaps; and whether it is desired in one
sector but unwanted in another. The light program may use a reflection-even
mode as a packet helper while the electric program requires the same branch to
avoid an additional long-range force. That is one shared constraint, not two
independent design choices.

The matrix also records the difference between an exact structural gate and a
bounded physical departure. Exact constituent conservation, a nonredundant
branch closure, bounded physical energy or branch-appropriate stability, the
absence of ghosts, the correct number of freely selectable light polarizations,
a legitimate force construction, and one shared normalization cannot be traded
against performance elsewhere. Leakage, resonance lifetime, weak dispersion,
extra radiation, species variation, return screening, and low-frequency drag
may instead be bounded over a declared regime. Detailed numerical budgets and
evidence paths remain in the capability matrix and derivation ledger rather
than being reproduced here.

\section{Shared fields, coefficients, speeds, sources, and constraints}

The common-medium test begins by identifying quantities that must have one
meaning everywhere. The first group belongs to the host material: the
background profiles \(n_0(w)\) and \(\chi_0(w)\), the selected thickness
\(H_{\rm br}\), the ordered-state shear law, the equation of state, and the
chosen conservative, relaxational, or explicitly mixed order dynamics. A sector
may not use a conservative free-energy Hessian while another silently treats the
same order variable as dissipative without the associated reservoir. Different
branch choices define different candidates because they change the phase space,
constraints, conserved quantities, force formalism, damping, entropy ledger,
and particle-longevity problem.

The second group is the complete constrained response. Where an equilibrium
reduction exists, the static calculation uses one physical Hessian
\(\mathcal H(\mathbf k)\). Dynamics uses one retarded operator
\(\mathcal O_R(\omega,\mathbf k)\), including every internal variable needed by
the selected branch. Their parity blocks, constraints, poles, branch cuts,
continua, norm signs, damping rates, profiles, and residues are shared data. The
transverse speed \(c_\gamma\), longitudinal speed \(c_L\), bulk compressional
speed \(c_s\), and any electric-sector speed or relaxation scale are therefore
outputs of one spectrum. They cannot be assigned separately to obtain light
confinement, Coulomb range, gravity propagation, and low drag.

The third group is the source hierarchy. The reflection-odd and
reflection-even source vectors, the oriented density and current, the leading
nonnegative electric coupling magnitude \(g_a\), the gravity source-side
reduction, and the normalization of outgoing transverse modes must retain one
parent definition. A provisional far-field source may leave its amplitude and
form factor symbolic, but its transformation law and continuity assumptions
must be explicit. The later throat program must derive those quantities. It may
select among source families permitted by Candidate V1; it may not alter the
background operator or introduce a new electric, magnetic, radiative, or
gravity coupling after seeing the result.

Boundary and ensemble data are shared as well. Reflection covariance requires a
specified transformation of the environmental data. An equilibrium pair, a
conservative periodic pair, and a driven dissipative pair do not automatically
share one interaction potential. Outgoing-wave conditions, fixed source or
fixed boundary-value ensembles, support-mode action or phase conditions, and
reservoir asymptotics must therefore be frozen with the candidate. Otherwise a
force sign or lifetime could be changed by changing the problem rather than the
medium.

These shared objects create direct cross-sector tensions. Increasing a
thickness-mode gap may screen an unwanted even electric force while weakening a
photon helper. A higher-gradient term may aid packet localization while adding
vacuum dispersion or a low-threshold branch. Giving the odd Coulomb carrier
ordinary inertia may produce a clean wave but also extra radiation; making it
relaxational may replace radiation with lag, friction, and reservoir cost.
Locking characteristic speeds may improve an emergent common cone while placing
a guided mode at a bulk continuum threshold. A distributed porous-brane
\(S_{\mathrm{leak}}\) return law that supplies the desired local gravity regime
changes the asymptotic reservoir seen by the slab, light, and throat. The purpose
of the shared-quantity map is to make these consequences unavoidable.

\section{Major cross-sector compatibility gates}

Seven architecture gates collect the largest known ways in which separately
reasonable requirements may conflict. They summarize the detailed companion
review rather than replacing it. Each gate begins as a present burden. An
adverse calculation normally rejects a branch or Candidate V1 first. It becomes
an architecture-level result only if the obstruction persists throughout the
admitted common-medium class or can be avoided only by prohibited retuning or
an unsupported new mechanism.

\begingroup
\small
\begin{longtable}{@{}L{0.12\textwidth}L{0.39\textwidth}L{0.41\textwidth}@{}}
\caption{Proposal-level summary of the seven cross-sector architecture gates.}
\label{tab:cross-sector-gates}\\
\toprule
\textbf{Gate} & \textbf{Common burden} & \textbf{Earliest decisive test and initial adverse scope} \\
\midrule
\endfirsthead
\toprule
\textbf{Gate} & \textbf{Common burden} & \textbf{Earliest decisive test and initial adverse scope} \\
\midrule
\endhead
1. Frozen candidate
& The broad ontology still permits more than one dynamical branch, operator
  basis, shear realization, higher-gradient completion, and reservoir closure.
  Results cannot be combined while sectors choose different members of that
  menu.
& Freeze Candidate V1, then solve its finite slab and complete constrained
  response. Failure initially rejects V1; it threatens the architecture only if
  no allowed frozen candidate can realize the contracts. \\
\addlinespace
2. Odd Coulomb dynamics
& A healthy reflection-odd branch with leading \(k^2\) static stiffness can
  produce a Coulomb tail, but that fact does not determine whether the branch
  propagates, relaxes, or is a genuine constrained variable.
& Derive the complete odd static and retarded blocks, constraints, poles,
  residues, radiation, passivity, and low-frequency friction. Static success
  does not pass if every completion creates a ghost, strong coupling,
  unacceptable radiation, drag, or independent electromagnetic normalizations. \\
\addlinespace
3. Speeds and confinement
& Relations among \(c_\gamma\), \(c_L\), \(c_s\), electric response, and slower
  branches affect guided-mode localization, continuum thresholds, gravity
  response, mode conversion, and the critical velocity of a moving throat.
& Compute the full spectrum, essential continua, source and interface overlaps,
  leakage widths, and directional all-branch critical velocity. Exact cone
  locking may survive through a derived decoupling; distinct speeds may survive
  only as an explicitly bounded preferred-medium analog. \\
\addlinespace
4. Universality and discreteness
& Distinct throat species must coexist with sufficiently universal unit electric
  coupling, positive and universal low-energy mass ratios, and a conserved
  oriented count or current. Continuous internal moduli may instead produce
  arbitrary charge, size, throughput, or mass.
& Continue solved throat families and measure source residues, response ratios,
  topology, and sensitivity to internal parameters. Failure initially rejects
  the proposed particle family or its charge and mass interpretation, not the
  already-derived background response. \\
\addlinespace
5. Photon binding and linearity
& The nonlinearity and dispersion strong enough to bind one finite packet also
  affect ordinary light speed, vacuum dispersion, polarization, wakes,
  collisions, even-sector forces, and throat support.
& From the same projected V1 coefficients, test boundedness, scaling, collapse,
  speed shift, helper radiation, and preliminary packet interactions. Failure
  of the shared binding window rejects the self-bound photon ontology before it
  rejects a viable linear guided-light sector. \\
\addlinespace
6. Localized drainage, distributed return, and reservoir closure
& Persistent one-way throat throughput and ordered-to-bulk conversion link
  throat longevity to energy, momentum, entropy, drag, gravity screening, and
  depletion of the common bulk reservoir. Compensating material return is a
  distinct process: bulk reorganizes into brane state across the porous brane
  through the positive bulk-to-brane contribution to \(S_{\mathrm{leak}}\).
& Iterate the local supported-throat problem with the global bulk-reservoir and
  distributed porous-brane return problem to a stable branch-consistent fixed
  point. Failure rejects the selected conservative, dissipative, or mixed
  throat branch; it threatens the architecture only if no branch can close. \\
\addlinespace
7. Contact and topology
& Cancellation of the leading odd electric disturbance, projected coincidence,
  four-dimensional core contact, topology change, and annihilation are distinct.
  Same-species partners must also be distinguished from opposite charges of
  different internal or composite species.
& Solve both orientations, continue two-core families through graph breakdown,
  and evolve conserving nonlinear collisions. Failure first rejects a simple
  annihilation mechanism or species map; the particle--antiparticle sector fails
  only if no conserving channel exists. \\
\bottomrule
\end{longtable}
\endgroup

The first three gates can be attacked substantially at the background and
linear-response stage. The photon gate begins there and continues through the
finite-amplitude light calculation. The universality, reservoir, and topology
gates cannot be finally closed without solved throats and global evolution, but
they constrain Candidate V1 before that work begins. The candidate must contain
an admissible conservation structure, nonlinear field content, branch-specific
reservoir, and topology capable of posing those tests. A later failure may reveal
that V1 lacks the required enforcing structure; it may not be repaired by adding
one only to the failed sector without declaring a new candidate version.

\section{The common viable-region test}

Let \(\lambda\) denote the independent continuous parameters of one
declared candidate architecture, with discrete branch, symmetry, operator, and
boundary choices carried as explicit labels rather than hidden coordinates. For
each mandatory requirement \(i\), let \(\Lambda_i\) be the subset on which its
exact gate closes or its bounded departure remains inside the declared
admissible set. The common contract-feasible region is

\[
\boxed{
\Lambda_{\rm contract}
=
\bigcap_{i\in\mathcal I_{\rm mandatory}}\Lambda_i.
}
\]

Equivalently, at the architecture level the principal burdens may be displayed
schematically as

\[
\begin{aligned}
\Lambda_{\rm architecture}
={}&
\Lambda_{\rm slab}
\cap\Lambda_{\rm guided\ light}
\cap\Lambda_{\rm Coulomb}
\cap\Lambda_{\rm electric\ dynamics}
\\
&\cap\Lambda_{\rm throat\ universality}
\cap\Lambda_{\rm photon\ packet}
\cap\Lambda_{\rm reservoir}
\cap\Lambda_{\rm topology}.
\end{aligned}
\]

Before throats are solved, some of these sets express necessary structural
conditions rather than completed existence results. A nonempty
\(\Lambda_{\rm contract}\) therefore authorizes selection of a realization
attempt; it does not yet prove that the background, throats, or reservoir fixed
point exist. Candidate V1 must convert the provisional intersection into a
realized one by producing the objects used to evaluate it.

Mandatory requirements are conjunctive. No weighted score permits a long photon
lifetime to compensate for failed charge conservation, a clean Coulomb tail to
compensate for a ghost, or an attractive gravity regime to compensate for an
external energy source disguised as an internal reservoir. Optional or stronger
capabilities may be reported separately, and a failure may reduce the model's
scope, but they cannot pay for a failed mandatory gate. Every point reported as
common must use the same parent version, branch, background, source conventions,
and shared parameter values.

The geometry of the surviving set is itself a scientific result. An empty set
means that the tested combination cannot be realized. An isolated point or
singular locus may indicate a special exact solution but provides no robustness
without a derived protection. A narrow band close to a continuum threshold,
pole collision, sign reversal, instability, collapse boundary, or vanishing
source residue is fragile. Several disconnected components represent distinct
candidate branches and must not be blended. A robust region contains a finite
neighborhood in which all hard gates remain closed and every bounded departure
retains a margin larger than the relevant numerical and model uncertainty.
Perturbing shared coefficients, boundary data, regulator choices, and retained
higher-order terms within their declared uncertainty should not immediately
reverse the decision.

The matrix should therefore report not only passing points but the limiting
requirements, distances to decision boundaries, uncertainty on those margins,
and which components of the intersection remain conditional on later throat or
reservoir work. A broad exploratory scan may locate promising regions or no-go
surfaces. The final claim, however, is the common region after one candidate and
its evaluation rules have been frozen.

\section{Freezing Candidate Parent System V1}

If the contract-feasible intersection is nonempty, the project selects one
complete realization for a full test. Candidate Parent System V1 is not the
entire architecture and is not a menu from which each sector may choose. It can
be represented schematically as

\[
\boxed{
\mathfrak{C}_1
=
\left(
\mathcal A_1,
\Phi_1,
\mathcal E_1,
\mathcal R_1,
\Sigma_1,
\lambda_1\in\Lambda_1,
\mathcal Q_1,
\mathcal O_1
\right),
}
\]

where the entries denote the parent arena; complete field content; equations,
constitutive laws, and operator basis; internal reservoir closure; symmetries
and constraints; independent parameter point or domain; boundary, source, and
ensemble rules; and operational maps used to extract observables. The explicit
freeze must state the maximum derivative order, any allowed nonlocality or
memory, time-reversal and reflection properties, the physical constrained
space, asymptotic and outgoing conditions, and the one-time calibration choices.

The field inventory includes every fundamental, retained effective, auxiliary,
constrained, and reservoir variable needed by the selected branch. The equation
set must specify whether the order sector is conservative, relaxational, or
mixed and must close the corresponding material, momentum, energy, and, where
applicable, entropy transfers. The operator-basis freeze is as important as the
numerical parameter freeze. A term omitted from V1 is absent for V1; it cannot
reappear later as a photon-only, electricity-only, or throat-only repair.

Candidate V1 may retain a declared common parameter domain when a family rather
than one number is being tested. It may not retain materially different field
contents or branch choices under one name. Alternative completions remain
separate candidates. Simplicity is a discriminator only after mandatory
capabilities are represented: omitting a required degree of freedom is not a
virtue, while adding speculative operators that no contract requires weakens the
meaning of later failure.

The freeze also identifies what remains to be derived. The stable thickness,
mode profiles, electric eigenvector, throat radius, source magnitudes,
universality, packet solution, distributed porous-brane
\(S_{\mathrm{leak}}\) return profile, and converged reservoir state are not
independent knobs merely because they are not yet known. They are outputs of V1.
Once obtained, they become inherited inputs for downstream calculations with
their regimes and uncertainty attached.

Freezing does not certify that V1 is correct. It makes the test definite. If a
later calculation exposes a missing assumption or an inconsistent operator, the
project may define V2. That revision must state the change and reopen every real
dependency. Results derived from V1 remain historical evidence about V1; they
cannot be combined with V2 outputs as though the parent system had never
changed.

\section{Stable brane, complete spectrum, and far-field rederivation}

Candidate V1 first faces the least expensive decisive realization test: can its
own equations produce the host medium assumed by the capability contracts? The
background solve must yield a finite ordered slab with selected thickness,
appropriate reflection properties, positive shear support, bounded or
branch-appropriate stable energy, and stable or sufficiently metastable
interfaces. The background profiles and asymptotic reservoir data must solve the
same V1 equations and boundary conditions that will be used downstream.

The next task is the complete constrained linear response about that slab. The
calculation must inventory every physical transverse, longitudinal, interface,
density, phase, order, bulk, electric, conversion, and reservoir branch. For an
equilibrium sector it derives the physical static Hessian and its parity blocks.
For dynamics it derives the full retarded operator. The output includes
constraints and auxiliary variables; eigenvectors or left and right response
modes where appropriate; static stiffnesses; poles, cuts, and continua; norm and
energy signs; characteristic speeds; gaps and screening lengths; damping and
leakage; normal profiles; source and interface overlaps; flux normalizations;
and observable residues. A sector-specific truncation is acceptable only after
this full operator shows that the omitted variables decouple in the stated
regime.

This one background calculation supplies the common inputs to several tests at
once. It must recover a healthy two-polarization guided transverse branch or
record the failure of the selected light branch. It must identify the
longitudinal and bulk thresholds that control helper radiation and drag. It must
show whether a reflection-odd source-coupled eigenbranch has the required
\(k^2\) static stiffness and whether its retarded completion is propagating,
relaxational, mixed, auxiliary, or constrained. It must classify the even
thickness and order response, every additional long-range force, and the modes
available for gravity response, throat support, radiation, and uniform-motion
emission.

The far-field contracts are then rerun on this actual background. Guided-light
profiles, speed, dispersion, leakage, and nonlinear projection coefficients are
computed from V1 rather than imported from a provisional action. Gravity uses
the V1 response and source-side variables with one frozen calibration and a
distributed porous-brane \(S_{\mathrm{leak}}\) return ensemble. Static
electricity uses the complete V1 odd and even blocks, the
physical boundary ensemble, and one source normalization. Magnetism and
acceleration radiation use the same oriented current, coupling, transverse mode
norms, and retarded response. The mass and drag audits use the same spectrum,
reservoir, and source conventions. Every unwanted odd, even, longitudinal,
bulk, and conversion channel remains in the accounting.

At this stage the localized source may still be a controlled placeholder, but it
is now a source class permitted by V1 and coupled to V1's actual normalized
modes. The later throat solve must replace its amplitude, profile,
conservation law, and universality assumptions. Thus a successful rederivation
establishes a realized \emph{common far-field medium}, not yet a realized
particle sector. A candidate that has a stable slab but requires different
parameter points for guided light, Coulomb response, acceptable dynamics, and
low drag fails the common-medium test just as surely as a candidate with no
stable slab.

A useful realized region may be written

\[
\boxed{
\Lambda_{\rm far}^{(1)}
=
\Lambda_1
\cap\Lambda_{\rm slab}^{(1)}
\cap\Lambda_{\rm spectrum}^{(1)}
\cap
\bigcap_{S\in\mathcal S_{\rm far}}
\Lambda_S^{(1)}.
}
\]

Only if this set is nonempty, with every mandatory result evaluated on the same
V1 background and parameter point or region, does the program advance to the
supported-throat and particle work in \cref{ch:throat-particle-program}.

\section{Decision outcomes of the common-medium phase}

The common-medium phase can end with several scientifically distinct outcomes.
An empty contract intersection means that the tested capability combination
does not justify a Candidate V1. It becomes an architecture failure only after
the admitted common-medium class has been examined broadly enough to show that
the obstruction is structural rather than peculiar to one realization.

A nonempty but isolated or fragile intersection justifies a focused realization
attempt, not a robust success claim. Candidate V1 may be frozen to determine
whether the apparent point survives the full background, spectrum, convergence,
and perturbation tests. Its proximity to every threshold must remain part of the
result. Disconnected viable components should be treated as separate candidates
rather than combined into a single flexible system.

A contract-feasible Candidate V1 can still fail because it has no selected
stable slab, an unhealthy constrained spectrum, a wrong force sign, an
unacceptable extra channel, incompatible characteristic speeds, or an empty
realized far-field region. That is a controlled candidate failure and is one of
the principal intended outputs of the program. It identifies which capability
combination could not be realized by the frozen equations.

A weaker sector may also survive after a stronger claim fails. A healthy guided
linear-light branch may remain even if every shared self-bound photon packet
fails. A locally inverse-square gravity analog may remain while a larger-scale
return behavior limits its scope. Such outcomes must be declared as partial
analogs, not averaged into either complete success or complete failure. Whether
a stronger capability is mandatory for V1 must be stated in the candidate's
scope before its decisive test.

The positive exit condition for this chapter is deliberately limited: one
versioned parent system produces a stable common background and complete healthy
response and rederives all mandatory light and far-field capabilities in one
shared region. That result authorizes the nonlinear throat program. It does not
establish particle existence, charge discreteness, universal mass response,
annihilation, or the global reservoir fixed point. Those deferred gates remain
capable of rejecting the particle branch or Candidate V1 later. The value of the
common-medium phase is that either outcome---a surviving common background or a
controlled incompatibility---has a precise meaning without sector-specific
retuning.
